\documentclass[11pt,a4paper]{article}
\usepackage[utf8]{inputenc}
\usepackage[T1]{fontenc}
\usepackage{lmodern}
\usepackage[english]{babel}
\usepackage{amsmath,amssymb}
\usepackage[margin=2.4cm]{geometry}
\usepackage{booktabs}
\usepackage{xcolor}
\usepackage{listings}
\usepackage{hyperref}
\newcommand{\norm}[1]{\left\lVert #1 \right\rVert}

\lstset{language=Python, basicstyle=\ttfamily\small, keywordstyle=\color{blue!70!black}\bfseries,
  commentstyle=\color{green!40!black}\itshape, stringstyle=\color{red!60!black}, frame=single,
  backgroundcolor=\color{gray!6}, columns=fullflexible, keepspaces=true, showstringspaces=false,
  upquote=true}

\title{Python and numpy for Assignment 4\\\large Step-by-step guide, grows with the project}
\author{Filippo Ferroni}
\date{Version 1.2, 25 October 2026 --- Lessons 1--2: functions, linear algebra, random data}

\begin{document}
\maketitle

\paragraph{How to use this guide.}
Each lesson matches one piece of the project. You learn only the Python you need, when you need it.
Each lesson ends with exercises; solutions to the project TODOs are added \emph{after} you have written them yourself.

\tableofcontents

%--------------------------------------------------------------------
\section{Project map}
\begin{center}
\begin{tabular}{@{}lll@{}}
\toprule
File & What it does & Lesson \\
\midrule
\texttt{src/models.py} & least squares, ridge, mean squared error & 1 (this one) \\
\texttt{src/data.py} & generates synthetic data, loads the real dataset, train/test split & 2 \\
\texttt{src/experiment.py} & loop over $d$: train and measure the errors & 3 \\
\texttt{src/plots.py} & plots with matplotlib & 4 \\
\bottomrule
\end{tabular}
\end{center}

%--------------------------------------------------------------------
\section{Running a Python file}
\begin{itemize}
  \item A \texttt{.py} file is text with instructions executed top to bottom.
  \item From the terminal, inside the repo folder:
\begin{lstlisting}
python src/models.py
\end{lstlisting}
        (on a Mac the command is sometimes \texttt{python3}).
  \item In VS Code: open the \texttt{smml-double-descent} folder, open the file, then use the ``Run'' button at the top right.
  \item Lines starting with \texttt{\#} are \textbf{comments}: Python ignores them, they are for the reader.
\end{itemize}

%--------------------------------------------------------------------
\section{Variables and \texttt{print}}
\begin{lstlisting}
m = 3              # an integer
alpha = 0.1        # a decimal number (use a dot, not a comma)
name = "ridge"     # a string (text)
print("alpha is", alpha)
\end{lstlisting}
The \texttt{=} sign \textbf{assigns} a value to a name; it is not a mathematical equality.

%--------------------------------------------------------------------
\section{Functions}
A function is a recipe: it takes inputs, does some computation, returns a result.
\begin{lstlisting}
def square(x):
    """Return x squared."""   # docstring: explains what it does
    result = x ** 2           # ** is the power operator
    return result

print(square(4))   # prints 16
\end{lstlisting}
\begin{itemize}
  \item \texttt{def} defines the function; the \textbf{arguments} (inputs) go in parentheses.
  \item The body is \textbf{indented by 4 spaces}. Indentation tells Python which lines belong to the function.
        If it is wrong: \texttt{IndentationError}.
  \item \texttt{return} gives back the result. Without \texttt{return}, the function returns \texttt{None} (``nothing'').
  \item \texttt{raise NotImplementedError} means ``not written yet'': you find it in the TODOs and replace it with your code.
\end{itemize}

\paragraph{The \texttt{if \_\_name\_\_ == "\_\_main\_\_":} block.}
The code below this line runs only when you launch the file directly (\texttt{python src/models.py}),
not when another file does \texttt{import} of \texttt{models.py}. We use it for quick checks.

%--------------------------------------------------------------------
\section{numpy: vectors and matrices}
\begin{lstlisting}
import numpy as np        # import numpy under the short name np

y = np.array([1.0, 2.0, 3.0])            # vector with 3 entries
X = np.array([[1.0, 0.0],
              [1.0, 1.0],
              [1.0, 2.0]])               # 3 x 2 matrix (3 rows, 2 columns)

print(X.shape)       # (3, 2)
print(X.shape[0])    # 3 = number of rows = m
print(X.shape[1])    # 2 = number of columns = d
print(X[0])          # first row: [1. 0.]   (counting starts at 0!)
print(X[2, 1])       # row 3, column 2: 2.0
\end{lstlisting}
\textbf{Careful:} Python counts from $0$. The first element is \texttt{[0]}, not \texttt{[1]}.

%--------------------------------------------------------------------
\section{Operations: element-wise vs matrix product}
\begin{center}
\begin{tabular}{@{}lll@{}}
\toprule
Maths & numpy & Notes \\
\midrule
$A + B$ & \texttt{A + B} & element-wise, same shapes \\
$\alpha A$ & \texttt{alpha * A} & multiplies every entry by $\alpha$ \\
$v_i^2$ for every $i$ & \texttt{v ** 2} & element-wise power \\
$AB$ (matrix product) & \texttt{A @ B} & \textbf{not} \texttt{A * B}! \\
$A^\top$ & \texttt{A.T} & transpose \\
$I_d$ & \texttt{np.eye(d)} & $d\times d$ identity \\
$\frac{1}{m}\sum_i v_i$ & \texttt{np.mean(v)} & mean of the entries \\
$\norm{w}$ & \texttt{np.linalg.norm(w)} & Euclidean norm \\
$A^{+}$ & \texttt{np.linalg.pinv(A)} & pseudo-inverse \\
solve $Aw = b$ & \texttt{np.linalg.solve(A, b)} & better than computing $A^{-1}$ \\
\bottomrule
\end{tabular}
\end{center}
\textbf{Classic mistake:} \texttt{A * B} multiplies entry by entry, it is \emph{not} the row-by-column product.
The formulas in the slides almost always need \texttt{@}.

\paragraph{The shape rule.}
\texttt{A @ B} works only if the number of columns of \texttt{A} equals the number of rows of \texttt{B}:
$(p\times q)\,@\,(q\times r) = (p\times r)$. Otherwise numpy raises an error containing the word \texttt{mismatch}.
\textbf{Before writing a line, check the shapes by hand.}

%--------------------------------------------------------------------
\section{From formula to code: the 4-step method}
\begin{enumerate}
  \item Write down the mathematical formula.
  \item Give a name to each piece (e.g.\ $A$, $b$).
  \item Translate each symbol using the table in the previous section.
  \item Return the result with \texttt{return}.
\end{enumerate}

\subsection{Worked example: ridge (TODO 2)}
Formula (slides 9, p.~5): $(\alpha I + X^\top X)\,w = X^\top y$, i.e.\ $A w = b$ with $A = \alpha I + X^\top X$, $b = X^\top y$.
\begin{lstlisting}
def fit_ridge(X, y, alpha):
    d = X.shape[1]                     # number of features (columns of X)
    A = alpha * np.eye(d) + X.T @ X    # d x d matrix
    b = X.T @ y                        # vector with d entries
    w = np.linalg.solve(A, b)          # solves A w = b
    return w
\end{lstlisting}
Shape check with $X$ of shape $m\times d$ and $y$ with $m$ entries:
\begin{itemize}
  \item \texttt{np.eye(d)}: $d\times d$; \quad \texttt{X.T @ X}: $(d\times m)(m\times d) = d\times d$; \quad sum: $d \times d$. \checkmark
  \item \texttt{X.T @ y}: $(d\times m)(m) = d$. \checkmark
  \item \texttt{solve}: from a $d\times d$ matrix and a $d$-vector it returns $w$ with $d$ entries. \checkmark
\end{itemize}

%--------------------------------------------------------------------
\section{Common errors and what they mean}
\begin{center}
\begin{tabular}{@{}p{5.2cm}p{9.3cm}@{}}
\toprule
Message & Likely cause \\
\midrule
\texttt{IndentationError} & wrong indentation (use 4 spaces, do not mix tabs and spaces) \\
\texttt{NameError: name 'np' is not defined} & missing \texttt{import numpy as np} at the top of the file \\
\texttt{NotImplementedError} & you left the \texttt{raise NotImplementedError} line in the TODO \\
\texttt{ValueError: matmul: ... mismatch} & incompatible shapes in a \texttt{@}: check the shapes \\
\texttt{LinAlgError: Singular matrix} & \texttt{solve} with a non-invertible matrix (e.g.\ ridge with $\alpha=0$ and $d>m$) \\
\texttt{ModuleNotFoundError: numpy} & numpy not installed: \texttt{pip install -r requirements.txt} \\
\bottomrule
\end{tabular}
\end{center}

%--------------------------------------------------------------------
\section{Lesson 1 exercises}
\begin{enumerate}
  \item With $X$ of shape $3\times 2$: why does \texttt{X @ X} fail while \texttt{X.T @ X} works? What is the shape of the result?
  \item What is the difference between \texttt{X.T @ X} and \texttt{X @ X.T}? What are their shapes with $m=3$, $d=2$?
        (Hint: one of them appears in the minimum-norm formula, theory Sec.~4.3.)
  \item What does \texttt{np.array([1.0, -2.0, 3.0]) ** 2} print? And \texttt{np.mean} of that result?
  \item \textbf{TODO 1} in \texttt{models.py}: translate $w = X^{+}y$ with the 4-step method.
  \item \textbf{TODO 3} in \texttt{models.py}: translate $\frac{1}{m}\norm{Xw - y}^2$. (Hint: first the error vector, then the mean of the squares.)
\end{enumerate}

\subsection*{Solutions (exercises 1--3)}
\begin{enumerate}
  \item \texttt{X @ X} needs columns of $X$ ($2$) $=$ rows of $X$ ($3$): false, hence the error.
        \texttt{X.T @ X} is $(2\times 3)(3\times 2) = 2\times 2$.
  \item \texttt{X.T @ X} is $d\times d = 2\times 2$ (normal equations, ridge); \texttt{X @ X.T} is $m\times m = 3\times 3$
        (it appears in $w^{+} = X^\top(XX^\top)^{-1}y$).
  \item \texttt{[1. 4. 9.]}; the mean is $14/3 \approx 4.67$.
\end{enumerate}
\subsection*{TODO solutions (written by you, checked on 25/10)}
\begin{lstlisting}
def fit_least_squares(X, y):
    w = np.linalg.pinv(X) @ y          # w = X^+ y, minimum norm for every d
    return w

def mse(X, y, w):
    y_hat = X @ w                      # predictions on the set (m values)
    error = np.mean((y_hat - y) ** 2)  # mean of the squared errors
    return error
\end{lstlisting}
Check outputs: A $=[1,1]$, B $=[1,10]$, C $\approx[1.0,\,0.01]$, D $\approx 2\cdot10^{-32}$ (numerical zero: the computer
works with about 16 significant digits, so such a tiny value is a zero).
\textbf{Lesson from check C:} on the same data the minimum-norm solution puts a weight of $10$, ridge $0.01$:
this is the mechanism that removes the peak, seen on your own code.

%====================================================================
\section{Lesson 2: generating data (\texttt{src/data.py})}
\paragraph{Why.} To see double descent we need data whose ``truth'' we know:
$y = w^{*\top}x + \text{noise}$, with $x \in \mathbb{R}^D$ ($D=500$ ``weak'' features), $m=100$ training examples
and $2000$ test examples. The model uses only the first $d$ features: growing $d$ from $1$ to $500$ crosses the threshold $d=m$.

\subsection{Reproducible random numbers: the seed}
\begin{lstlisting}
rng = np.random.default_rng(42)   # generator with seed 42
a = rng.normal(size=3)            # 3 numbers from N(0,1)
\end{lstlisting}
\begin{itemize}
  \item \texttt{rng} is a ``generator'' object. The same seed \textbf{always} gives the same sequence: reproducible experiments
        (an assignment requirement).
  \item \texttt{rng.normal(size=...)} draws from a standard normal $N(0,1)$. \texttt{size} sets the shape:
        \texttt{size=5} gives a vector of 5 values, \texttt{size=(100, 500)} a $100\times 500$ matrix (a \textbf{tuple} in round brackets).
  \item Noise with standard deviation $\sigma$: \texttt{sigma * rng.normal(size=m)}, because if $Z\sim N(0,1)$ then $\sigma Z \sim N(0,\sigma^2)$.
\end{itemize}

\subsection{Normalizing a vector}
\texttt{w\_star / np.linalg.norm(w\_star)} divides every entry by the norm: the result has norm $1$.
So the signal $w^{*\top}x$ has variance $1$ and the signal-to-noise ratio is controlled by $\sigma$ alone.

\subsection{Slicing: taking a piece of a matrix}
\begin{center}
\begin{tabular}{@{}ll@{}}
\toprule
Code & Meaning \\
\midrule
\texttt{X[:, :d]} & all rows, columns $0$ to $d-1$ (the first $d$ features) \\
\texttt{X[:10]} & the first 10 rows \\
\texttt{X[:, 0]} & the first column (a vector) \\
\texttt{X[5:, 2:4]} & rows from 5 on, columns 2 and 3 \\
\bottomrule
\end{tabular}
\end{center}
Rule: \texttt{a:b} includes \texttt{a} and \textbf{excludes} \texttt{b}; a bare \texttt{:} means ``everything''.

\subsection{Returning several values}
\begin{lstlisting}
return X_train, y_train, X_test, y_test, w_star        # inside the function
X_tr, y_tr, X_te, y_te, w = make_synthetic(100, 2000, 500, 0.5, 0)   # caller
\end{lstlisting}
Python packs the values into a tuple and ``unpacks'' them in the order you wrote them.

\subsection{Why check C gives about 0.25}
With $\sigma = 0.5$ even the true $w^*$ is off by $\sigma^2 = 0.25$ on average: this is the \textbf{Bayes risk}
for the square loss (statRisk p.~1--2). No model can do better on the test set: in the plots it is the horizontal reference line.
The exact value ($0.256$) fluctuates a little because the test set is finite.

\subsection*{Lesson 2 exercises}
\begin{enumerate}
  \item What is the shape of \texttt{rng.normal(size=(3, 4))}? And of \texttt{X[:, :2]} if $X$ is $3\times4$?
  \item Why do we draw the test set with the \emph{same} \texttt{rng} but \emph{after} the training set? What does it guarantee?
  \item With \texttt{noise\_std=1.0}, what should check C approach?
\end{enumerate}
\subsection*{Solutions}
\begin{enumerate}
  \item $(3,4)$; $(3,2)$.
  \item The draws are independent: no example is shared between train and test (no data leakage), and the seed keeps everything reproducible.
  \item $1$ ($=\sigma^2$).
\end{enumerate}
\subsection*{TODO solutions (written by you, checked on 25/10)}
\begin{lstlisting}
X_train = rng.normal(size=(m, D))                             # TODO 1
y_train = X_train @ w_star + noise_std * rng.normal(size=m)  # TODO 2
return X[:, :d]                                               # TODO 3
\end{lstlisting}
Output: A correct shapes, B $(100,30)$, C $0.2564$, D \texttt{True}.
\textbf{Clean-up:} the \texttt{raise NotImplementedError} line left after the \texttt{return} never runs (dead code): delete it.

\end{document}
